\documentclass[12pt]{article}
\usepackage{amssymb}
\usepackage[fleqn]{mathtools}
\setlength{\mathindent}{0pt}%stânga

\usepackage{authblk}%pachet pentru introducerea afilierii
\usepackage[]{geometry}
\usepackage{graphicx}%rotește ecuații
\usepackage{amsthm}
\usepackage{amsmath, amssymb}
\usepackage{enumitem}
\usepackage{babel}
\usepackage{cancel}

\newenvironment{rightdisplay}
  {\begin{flushright}\(\displaystyle}
  {\)\end{flushright}} %ecuație aliniată cu marginea dreaptă

\theoremstyle{definition}
\newtheorem{defi}{Defini\c tia}
\newtheorem{teo}{Teorema}
\newtheorem{prop}{Propozi\c tia}
\newtheorem{lm}{Lema}
\newtheorem{cor}{Corolarul}
\newtheorem{obs}{Observa\c tia}
\newtheorem{ex}{Exerci\c tiul}
\newtheorem{exemp}{Exemplul}
\renewcommand{\proofname}{Demonstra\c tie}
%\renewcommand{\bibname}{Bibliografie}
\renewcommand{\contentsname}{Cuprins}
%\renewcommand{\chaptername}{Capitolul}
\renewcommand{\figurename}{Fig.}

\title{Generalizarea Șirurilor: Fibonacci, Lucas, Pell}
\author{\enspace}
\affil{\enspace}
\date{\enspace}

\begin{document}
	\section*{Șiruri}
		\subsection*{Fibonacci}
		\[\begin{aligned}
			& F_{n+2} = F_{n+1}+F_{n}, F_1=1, F_2=1 \\
			& 1,1,2,3,5,8,... \\
			& \lim\limits_{n\to \infty} \frac{F_{n+1}}{F_n}= \frac{1+\sqrt{5}}{2}
		\end{aligned}\]
		\subsubsection*{Termenul general}
		\[\begin{aligned}			
			& F_{n+2} = F_{n+1} + F_{n}, F_1 = 1, F_2 = 1\\
			& F_{n+2} - F_{n+1} - F_{n} = 0\\
			& \text{Ecuație caracteristică: } x^2 -x -1 = 0\\
			&\Delta = 5 => x_{1,2} = \frac{1 \pm \sqrt{5}}{2}\\
			&F_n = A\left(\frac{1+\sqrt{5}}{2}\right)^n
    				+ B\left(\frac{1-\sqrt{5}}{2}\right)^n\\
    			& F_n=\frac{1}{\sqrt{5}} \left[ \left(\frac{1+\sqrt{5}}{2}\right)^n-\left(\frac{1-\sqrt{5}}{2}\right)^n \right]
		\end{aligned}\]
		\subsection*{Lucas}
		\[\begin{aligned}
			&L_{n+2}=L_{n+1}+L_n, F_1=1, F_2=3\\
			&1,3,4,7,11,...
		\end{aligned}\]
		\subsubsection*{Termenul general}
		\[\begin{aligned}
			&L_{n+2}-L_{n+1}-L_n = 0\\
			&\text{Ecuație caracteristică: }x^2-x-1=0 \\
		&\Delta = 5 \Rightarrow x_{1,2} = \frac{1 \pm \sqrt{5}}{2}
		\Rightarrow L_n = A\left(\frac{1+\sqrt{5}}{2}\right)^n
						+ B\left(\frac{1-\sqrt{5}}{2}\right)^n
		\end{aligned}\]

		\[\begin{aligned}
			&\left.
				\begin{aligned}
					&L_0=A+B=2\\
					&L_n= A\left(\frac{1+\sqrt{5}}{2}\right)
       						+ B\left(\frac{1-\sqrt{5}}{2}\right)
				\end{aligned}
			\right\}
			\Rightarrow
			A\left(\frac{1+\sqrt{5}}{2}\right)+(1-\sqrt{5}) - A\left(\frac{1+\sqrt{5}}{2}\right)=1\\
			&A\sqrt{5}=\sqrt{5} \Rightarrow A=1 \Rightarrow B=1
		\end{aligned}\]
		
		\subsection*{Pell}
		\[\begin{aligned}
			&P_{n+2}=2P_{n+1}+P_n, F_1=1, F_2=2\\
			&1,2,5,12,29,\cdots
		\end{aligned}\]
		\subsubsection*{Termenul general}
		\[\begin{aligned}
			&P_{n+2}-2P_{n+1}-P_n = 0\\
			&\text{Ecuație caracteristică: }x^2-2x-1=0 \\
		&\Delta = 8 \Rightarrow x_{1,2} = \frac{2 \pm 2\sqrt{2}}{2}=1\pm\sqrt{2}
			\ \ \ \begin{aligned}
				&P_2 = 2 \cdot P_1 + P_0\\
				&2 = 2 + P_0 \Rightarrow P_0=0
			\end{aligned}
		\end{aligned}\]

		\[\begin{aligned}
			&\left\{
				\begin{aligned}
					&P_0=A+B=0\\
					&P_n= A(1+\sqrt{2})+B(1+\sqrt{2})=1 
				\end{aligned}
			\right.
			\Rightarrow
			A(1+\sqrt{2})-A(1-\sqrt{2})=1\\
			&2A\sqrt{2}=1 \Rightarrow A=\frac{1}{2\sqrt{2}} \Rightarrow B=-\frac{1}{2\sqrt{2}}
			&
		\end{aligned}\]
		\section*{Forma generală}
			\subsection*{Fibonacci}
				\[\begin{aligned}
				A\ &=\begin{pmatrix}
					1	&	1	\\
					0	&	1	
				\end{pmatrix}\\
				A^2&= \begin{pmatrix}
					1	&	1	\\
					0	&	1	
				\end{pmatrix} \begin{pmatrix}
					1	&	1	\\
					0	&	1	
				\end{pmatrix}=\begin{pmatrix}
					1	&	2	\\
					0	&	1	
				\end{pmatrix}\\
				A^3&=\begin{pmatrix}
					1	&	2	\\
					0	&	1	
				\end{pmatrix} \begin{pmatrix}
					1	&	1	\\
					0	&	1	
				\end{pmatrix}=\begin{pmatrix}
					1	&	3	\\
					0	&	1	
				\end{pmatrix}\\
				&\cdots\\
				A^n&=\begin{pmatrix}
					1	&	n	\\
					0	&	1	
				\end{pmatrix}\\
				\end{aligned}\]
				
				\[\begin{aligned}
				A\ &=\begin{pmatrix}
					1	&	1	\\
					1	&	0	
				\end{pmatrix}\\
				A^2&= \begin{pmatrix}
					1	&	1	\\
					1	&	0	
				\end{pmatrix} \begin{pmatrix}
					1	&	1	\\
					1	&	0	
				\end{pmatrix}=\begin{pmatrix}
					2	&	1	\\
					1	&	1	
				\end{pmatrix}  \\
				A^3&=\begin{pmatrix}
					2	&	1	\\
					1	&	1	
				\end{pmatrix} \begin{pmatrix}
					1	&	1	\\
					1	&	0	
				\end{pmatrix}=\begin{pmatrix}
					3	&	2	\\
					2	&	1	
				\end{pmatrix} \ \ \ \ \ \ \ \ \ \ \ \ \ \ \ \ \ \ \ \ 
				 							A^n = \begin{pmatrix}
													F_{n+1}	& F_{n}\\
													F_n		&	F_{n-1}
											\end{pmatrix}\ \ \ ^{(\boldsymbol{*})}\\
				A^4&=\begin{pmatrix}
					3	&	2	\\
					2	&	1	
				\end{pmatrix} \begin{pmatrix}
					1	&	1	\\
					1	&	0	
				\end{pmatrix}=\begin{pmatrix}
					5	&	3	\\
					3	&	2	
				\end{pmatrix}\\
				A^5&=\begin{pmatrix}
					5	&	3	\\
					3	&	2	
				\end{pmatrix} \begin{pmatrix}
					1	&	1	\\
					1	&	0	
				\end{pmatrix}=\begin{pmatrix}
					8	&	5	\\
					5	&	3	
				\end{pmatrix}\\
				&\cdots\\
				A^n&=\begin{pmatrix}
					F_n + F_{n-1}	&	F_n	\\
					F_n	&	F_{n-1}	
				\end{pmatrix}=\begin{pmatrix}
				F_n & F_n\\
				F_n & 0
				\end{pmatrix}+\begin{pmatrix}
				F_{n-1} & 0\\
				0		& F_{n-1}
				\end{pmatrix} \Rightarrow A^n = F_n \begin{pmatrix}
				1	&	1\\
				1	&	0
				\end{pmatrix}+F_{n-1} \begin{pmatrix}
				1	&	0\\
				0	&	1
				\end{pmatrix}
				\end{aligned}\]
				
	\[\begin{aligned}
	A^n=F_n \cdot A = F_n \cdot A + F_{n-1} \cdot I_2,\  n \ge 1
	\end{aligned}\]
	\[ \text{Din} (\boldsymbol{*}) \Rightarrow  \left\{\begin{aligned} &\det(A^n) = \begin{vmatrix}
	F_{n+1}	& F_{n}\\
	F_{n}	& F_{n=1}
	\end{vmatrix}= F_{n+1} \cdot F_{n-1} - {F_n}^2 = (\det{A})^n \\
	&\det{A} = \begin{vmatrix}
	1&1\\
	1&0
	\end{vmatrix}=-1 \end{aligned}
	\right. \Rightarrow\] \[ \Rightarrow \det{A^n} = F_{n+1} \cdot F_{n-1} - {F_n}^2 = (-1)^n \]
	\begin{obs} 
	\[\begin{aligned}
	\begin{pmatrix}
	F_{n+2}\\
	F_{n+1}
	\end{pmatrix} &= \begin{pmatrix}
	1&1\\
	1&0
	\end{pmatrix} \begin{pmatrix}
	F_{n+1}\\ F_n
\end{pmatrix}\\	
	\begin{pmatrix}
	F_{n+1}\\
	F_{n}
	\end{pmatrix} &= \begin{pmatrix}
	1&1\\
	1&0
	\end{pmatrix} \begin{pmatrix}
	F_{n}\\ F_{n-1}
	\end{pmatrix}\\
	 &\cdots \\
	\begin{pmatrix}
	F_{2}\\
	F_{1}
	\end{pmatrix}\  &= \begin{pmatrix}
	1&1\\
	1&0
	\end{pmatrix} \begin{pmatrix}
	F_{1}\\ F_0
	\end{pmatrix} \xRightarrow{+} \begin{pmatrix}
	F_{n+1}\\
	F_{n}
	\end{pmatrix} = \begin{pmatrix}
	1&1\\
	1&0
	\end{pmatrix}^n \begin{pmatrix}
	F_1\\ F_0 
	\end{pmatrix}=A^n \begin{pmatrix}
	F1 \\ F0
	\end{pmatrix}
	\end{aligned}\]
	\end{obs}
\subsection*{Pell}
\[\begin{aligned}
A \ &=\begin{pmatrix}
2&1\\
1&0
\end{pmatrix}\\
A^2 &= \begin{pmatrix}
2&1\\
1&0
\end{pmatrix}\begin{pmatrix}
2&1\\
1&0
\end{pmatrix}=\begin{pmatrix}
5&2\\
2&1
\end{pmatrix}\\
A^3 &=  \begin{pmatrix}
5&2\\
2&1
\end{pmatrix}\begin{pmatrix}
2&1\\
1&0
\end{pmatrix}=\begin{pmatrix}
12&5\\
5&2
\end{pmatrix}\\
A^3 &= \begin{pmatrix}
12&5\\
5&2
\end{pmatrix}\begin{pmatrix}
2&1\\
1&0
\end{pmatrix}=\begin{pmatrix}
29&12\\
12&5
\end{pmatrix}\\
&\cdots
A^n &= \begin{pmatrix}
P_{n+1}&P_n\\
P_n&P_{n-1}, \ \ n \ge 1
\end{pmatrix} 
\end{aligned}\]
$$\det{A^n} = \begin{vmatrix}
P_{n+1}&P_n\\
p_n & P_{n-1}
\end{vmatrix}= P_{n+1}\cdot P_{n-1} - {P_n}^2 = (-1)^n$$
\[\begin{aligned}
A^n &= \begin{pmatrix}
2P_n + P_{n-1}& P_n\\
P_n&P_{n-1}
\end{pmatrix}=Pn\begin{pmatrix}
2&1\\
1&0
\end{pmatrix}+\begin{pmatrix}
P_{n-1}&0\\
0&P_{n-1}
\end{pmatrix}\\
&= P_n A+P_{n-1}I_2 
\end{aligned}\]
\subsection*{Lucas}
\[
\begin{aligned}
&\overset{\scalebox{1.3}{\ \ B}}{
  \overset{\scalebox{2}[1]{\ \rotatebox{90}{\}}}}{
    \begin{pmatrix}
      1 & 2\\
      2 & -1
    \end{pmatrix}
  }
}
\cdot
\overset{\scalebox{1.3}{\ \ A}}{
  \overset{\scalebox{1.6}[1]{\ \rotatebox{90}{\}}}}{
    \begin{pmatrix}
      1 & 1\\
      1 & 0
    \end{pmatrix}
  }
}
=
\begin{pmatrix}
  3 & 1\\
  1 & 2
\end{pmatrix}\\
&\begin{pmatrix}
  3 & 1\\
  1 & 2
\end{pmatrix}
\begin{pmatrix}
  1 & 1\\
  1 & 0
\end{pmatrix}
=
\begin{pmatrix}
  4 & 3\\
  3 & 1
\end{pmatrix}\\
&\begin{pmatrix}
  4 & 3\\
  3 & 1
\end{pmatrix}
\begin{pmatrix}
  1 & 1\\
  1 & 0
\end{pmatrix}
=
\begin{pmatrix}
  7 & 4\\
  4 & 3
\end{pmatrix}
\end{aligned}
\]
\[\begin{aligned}
&B \cdot A^{n-1} = \begin{pmatrix}
L_{n+1}&L_n\\
L_n&L_{n-1}
\end{pmatrix}, n \geq 1\\
&B\cdot\begin{pmatrix}
F_n&F_{n-1}\\
F_{n-1}&F_{n-2}
\end{pmatrix}= \begin{pmatrix}
1&2\\
2&-1
\end{pmatrix}\cdot\begin{pmatrix}
F_n&F_{n-1}\\
F_{n-1}&F_{n-2}
\end{pmatrix}=\begin{pmatrix}
F_{n}+2F_{n-1}&F_{n-1}+2F_{n-2}\\
2F_n-F_{n-1}&2F_{n-1}-F_{n-2}
\end{pmatrix}=\\
&L_{n+1}=F_{n}+2F{n-1}\\
&\boxed{L_n=F_{n-1}+2F_{n-2}}
\end{aligned}\]

\subsection*{Generalizare}
\[\left. \begin{aligned}
&x_{n+2} = k x_{n+1}+x_n & n=0\\
&x_1=1 & x_2=k x_1 + x_0\\
&x_2=k & \cancel{k} = \cancel{k} + x_0 
\end{aligned}\right\} A = \begin{pmatrix}
k&1\\
1&0
\end{pmatrix} \ \ \ \begin{aligned} &k=1, &A^n=\begin{pmatrix}
x_{n+1}&x_n\\
x_n&x_{n-1}
\end{pmatrix} \ \text{(Fib.)}\\
&&x_{n+2}=x_{n+1}+x_n \  \text{(Fib.)}\\
&k=2, &x_{n+2}=2x_{n+1}+x_n \ \text{(Pell)}
\end{aligned}\]

\subsection*{Identități legate de șiruri}
\[\begin{aligned}
&F_{n+2}=F_{n+1}+F_n \ |\ \cdot F_{n+1} \ \ F_1=1 \ F_2=1\\
&F_{n+1}F_{n+2}=F^2_{n+1} + F_n F_{n+1} \Rightarrow
\end{aligned}\]
\[\begin{aligned}
\Rightarrow F^2_{n+1}&=F_{n+1}F_{n+2} - F_n F_{n+1}\\
n\rightarrow n-1: F^2_{n}&=F_{n}F_{n+1} - F_{n-1} F_{n}\\
n\rightarrow n-2: F^2_{n-1}&=F_{n-1}F_{n} - F_{n-2} F_{n-1}\\
&\cdots \\
n=0 \rightarrow F^2_{1} &= F_1F_2-F_0F_1\\
\rule{3cm}{0.4pt}&\rule{4cm}{0.4pt}\ _+\\
\end{aligned}\]
\[\begin{aligned}
&\boxed{F^2_1+F^2_2+\cdots+F^2_n=F_n\cdot F_{n+1}}\\
\end{aligned}\]
\[
\begin{aligned}
&x_{n+2}=kx_{n+1}+x_n, \ \ x_1=1, \ x_2=k=kx_1+x_0 \Rightarrow x_0=0\\
&\text{generalizarea pentru șirul lui Fibonacci k=1}
\end{aligned}\]
\[\begin{aligned}
x_{n+2}x_{n+1} &= kx^2_{n+1} +x_nx_{n+1}\\
\Leftrightarrow kx^2_{n+1} &= x_{n+2}x_{n+1} - \cancel{x_nx_{n+1}}\\
n \rightarrow n-1\ kx^2_{n} &= \cancel{x_{n+1}x_{n}} - \cancel{x_{n-1}x_{n}}\\
n \rightarrow n-2\ kx^2_{n-1} &= \cancel{x_{n}x_{n-1}} - \cancel{x_{n-2}x_{n-1}}\\
&\cdots\\
n = 0\ kx^2_{1} &= \cancel{x_{2}x_{1}} - x_{0}x_{1}\ _{(=0)}\\
\rule{3cm}{0.4pt}&\rule{4cm}{0.4pt}\ _+
\end{aligned}
\]
\[\begin{aligned}
&\boxed{k(x^2_1 + x^2_2 + x^2_3 + \cdots + x^2_n)=x_n x_{n+1}} 
\end{aligned}\]
Dacă $k=2 \Rightarrow x_n = P_n $ (șirul lui Pelle), și:\\
$2(P^2_1 + P^2_2 + P^2_3 + \cdots + P^2_n)=P_n P_{n+1}$
\[
\sum_{k=1}^{n} \frac{F_k}{F_{k+1} F_{k+2}} = 1 - \frac{1}{F_{n+2}},\ n \geq 1
\]
\[
\frac{F_1}{F_2F_3}+\frac{F_2}{F_3F_4}+\cdots+\frac{F_n}{F_{n+1}F_{n+2}}= 1 - \frac{1}{F_{n+2}}
\]
\[
	\left( \frac{F_k}{F_{k+1}F_{k+2}} = \frac{F_{k+2}-F_{k+1}}{F_{k+1}F_{k+2}} =\frac{F_{k+2}}{F_{k+1}F_{k+2}}-\frac{F_{k+1}}{F_{k+1}F_{k+2}} = \frac{1}{F_{k+1}} - \frac{1}{F_{k+2}} \right)
\]
\[
\begin{aligned}
	\Rightarrow \sum_{k=1}^{n} \frac{F_k}{F_{k+1} F_{k+2}} & = \frac{1}{F_2} -\cancel{\frac{1}{F_3}} + \cancel{\frac{1}{F_3}} - \cancel{\frac{1}{F_4}} + \cancel{\frac{1}{F_4}} -\cancel{\frac{1}{F_5}} + \cdots + \cancel{\frac{1}{F_{n+1}}} - \frac{1}{F_{n+2}} \\ 
	&= \frac{1}{F_2} - \frac{1}{F_{n+2}}= 1 - \frac{1}{F_{n+2}}
\end{aligned}
\]
\begin{obs}
	\[ \sum_{k=1}^{n} \frac{F_k}{F_{k+1} F_{k+2}}  < 1\]
\end{obs}

\subsection*{General}
\[
 \sum_{i=1}^{n} \frac{x_i}{x_{i+1} x_{i+2}}=\frac{x_{i+2}-kx_{i+2}}{x_{i+2}}=\frac{1}{x_{i+1}}-\frac{k}{x_{i+2}}=kx_{i+1} + x_0 
\]
\[
	\begin{aligned}
		\frac{x_{n}}{x_{n+1}x_{n+2}}&=\frac{1}{x_{n+1}}-\frac{k}{x_{n+2}}\ |\cdot k^{n-1}\ \Leftrightarrow  &k^{n-1}\frac{x_{n}}{x_{n+1}x_{n+2}}&=\cancel{\frac{k^{n-1}}{x_{n+1}}}-\frac{k^n}{x_{n+2}}\\
		\frac{x_{n-1}}{x_{n}x_{n+1}}&=\frac{1}{x_{n}}-\frac{k}{x_{n+1}}\ |\cdot k^{n} \ \ \ \ \ \ \ \cdots & &\cdots \\
		&\cdots \ \ \ \ \ \ \ \ \ \ \ \ \ \ \ \ \ \ \ \ \ \ \ \ \ \ \ \ \cdots & k^2\frac{x_{3}}{x_{4}x_{5}}&=\cancel{\frac{k^2}{x_{4}}}-\cancel{\frac{k^3}{x_{5}}}\\
		\frac{x_{2}}{x_{3}x_{4}}&=\frac{1}{x_{3}}-\frac{k}{x_{4}}\ |\cdot k^1 \ \ \ \ \ \ \ \ \ \   \Leftrightarrow &k\frac{x_{2}}{x_{3}x_{4}}&=\cancel{\frac{k}{x_{3}}}-\cancel{\frac{k^2}{x_{4}}}\\
		\frac{x_{1}}{x_{2}x_{3}}&=\frac{1}{x_{2}}-\frac{k}{x_{3}}\ |\cdot k^0 \ \ \ \ \ \ \ \ \ \   \Leftrightarrow &\frac{x_{1}}{x_{2}x_{3}}&=\frac{1}{x_{2}}-\cancel{\frac{k}{x_{3}}}\\
		&&\rule{2cm}{0.4pt}&\rule{3cm}{0.4pt}\ _+
	\end{aligned}
\]
\[
\ \ \ \ \ \ \ \ \ \ \ \ \ \ \ \ \ \ \ \ \ \ \ \ \ \ \ \ \ \ \ \ \ \ \ \ \ \ \ \ \ \ \ \ \ \ \ \ \ \ \ \ \ \ \ \boxed{\sum_{i=1}^{n} k^{i-1}\frac{x_i}{x_{i+1} x_{i+2}}= 1-\frac{k^n}{x_{n+2}}}
\]
\[\ \ \ \ \ \ \ \ \ \ \ \ \ \ \ \ \ \ \ \ \ \ \ \ \ \ \ \ \ \ \ \ \ \ \ \ \ \ \ \ \ \ \ \ \ \ \ \ \ \ \ \ \ \ \ \  k=1 \rightarrow \text{Fib.}\ \  k=2\rightarrow\text{Pell} ??????\]

\subsection*{Cuaternari (Hamilton)}
\[\begin{aligned}
&z=a+ib\ \ \ a,b \in \mathbb{R},\ \ i^2=-1\\
&q=q_0+q_1i+q_2j+q_3k\\
&i^2 = j^2 = k^2 = -1\ \ \ ijk=-1\\
&ij = k =-ji\ \ \ \ jk=i=-kj \ \ \ \ ki=j=-ik 
\end{aligned}\]
Horadam introduce al n-lea cuaternion Fibonacci:
\[Q_n=F_n + F_{n+1}i + F_{n+2}j + F_{n+3}k\]
\subsection*{Teorema lui B\'{e}zoat}
\begin{teo}
Pentru orice $a,b \in \mathbb{Z}$ există $x,y\in \mathbb{Z}$ astfel încât $ax+by=(a,b)$ \\(i.e. CMMDC.)
\end{teo}
\begin{obs}
Dacă $(a,b) = 1$ i.e. $a$ și $b$ sunt prime între ele (coprime) atunci $(\exists) x,y \in \mathbb{Z}$ astfel încât $ax+by=1$   
\end{obs}
Descompuși în baze factori primi $a$ și $b$:
\[\left.\begin{aligned} 
&a = p^{\alpha_1}_1 \cdot\  \cdots\ \cdot p^{\alpha_k}_k \cdot u^{\alpha_{k+1}}_{k+1} \cdot \ \cdots \ \cdot u^{s}_{s}\\
&b = p^{\beta_1}_1 \cdot\  \cdots\ \cdot p^{\beta_k}_k \cdot v^{\beta_{k+1}}_{k+1} \cdot \ \cdots \ \cdot v^{t}_{t}
\end{aligned} \right\} \Rightarrow (a,b)=\prod_{i=1}^{k} p^{\min\{\alpha_i,\beta_i\}}_i
\]
Exemplu:
\[\begin{aligned}
&24 = 2^3 3 & \ \\
&\underline{36 = 2^2 3^2} &\  24x+36y&=12\\
&Cmmdc=2^2 3^1 = 12 \ \ \  &\  24(-1)+36(1)&=12
\end{aligned}\]

\subsubsection*{Metoda Blankinship}
Se aduce matricea $A$ formată dintr-o coloană cu numerele pentru care se caută CMMDC-ul, extinsă cu o matrice de forma $I_n$ și se aduce prin transformări elementare la una dintre cele 2 forme de mai jos (ex. cu 2 numere, $n=2$):
\[\begin{pmatrix}
a&1&0\\
b&0&1
\end{pmatrix} \sim \begin{pmatrix}
d&x&y\\
0&x_1&y_1
\end{pmatrix} \text{sau} \begin{pmatrix}
0&x_1&y_1\\
d&x&y
\end{pmatrix} \text{unde } d=(a,b), \ a,b,d,x,y,x_1,y_1\in\mathbb{Z}
\]
Analogie cu sisteme de ecuații:\[
\left\{\begin{aligned}
x+y&=2\\
3x+2y&=5
\end{aligned} \right. \text{ analog cu } \begin{pmatrix}
1&3&2\\
1&2&5
\end{pmatrix}
\]
Reamintire transformări elementare asupra liniilor unei matrice:
\begin{enumerate}[label={$T_{\arabic*}: \ $}]
  \item Înmulțirea unei linii printr-un număr real nenul.
  \item Schimbarea a 2 linii între ele.
  \item Înlocuirea unei linii prin suma dintre această linie și o altră linie înmulțită cu un număr real.
\end{enumerate}

\subsection*{Șirul lui Mersenne}

\[\begin{aligned}
&M_{n+2}=3M_{n+1}-2M_n \ \ \ \ M_1=1, \ M_2=3 \\
&M_{n+2}-3M_{n+1}+2M_n=0\\
&\text{Atașăm ecuația caracteristică: }x^2-3x+2=0\\
&x_1=2 \ \ x_2=1\\
&\\
&X_n=A{x_1}^n+B{x_2}^n, \ M_n=A\cdot 2^n +B\\
	&\left\{ \begin{aligned}
		&M_2=2=4A+B\\
		&M_1=1=2A+B
	\end{aligned} \right. \Rightarrow \left. \begin{aligned}
									&A=1\\
									&B=-1
									\end{aligned} \right\} \Rightarrow 2^n-1
\end{aligned}\]
\[\begin{aligned}
&1,3,7,15,31,...
\end{aligned}\]

\[\begin{aligned}
& 2018:\ M_{82589933}\\
& 2024:\  M_{136279841}
\end{aligned}\]

...
%\section*{Teme:}
%\subsection*{Metoda Blankenship}
%\[
%\begin{pmatrix}
%4932 & 1 & 0 & 0\\
%3324 & 0 & 1 & 0\\
%2013 & 0 & 0 & 1
%\end{pmatrix} =

%\]

%	\maketitle
%	\tableofcontents

%\begin{thebibliography}{3}
%	\bibitem{Angel} \textbf{Angel,A.R., Porter,S.R.}, \textit{A Survey of Mathematics with Applica-
%		tions,}, New York, 1997
%	\bibitem{Batschelet} \textbf{Batschelet,E.,} \textit{Introduction to Mathematics for Life Scientists,} Springer Verlag, Berlin-Heidelberg-New York, 1971.
%	\bibitem{Cret} \textbf{Creț,F.,} \textit{Elemente de Modelare și Matematici Speciale}, Editura MIRTON, Timișoara, 2000
%	
%\end{thebibliography}
%
\end{document}